% Сборка: pdflatex lab03.tex
\documentclass[a4paper,12pt]{article}

\usepackage{cmap}
\usepackage[T1,T2A]{fontenc}
\usepackage[utf8]{inputenc}
\usepackage[english,russian]{babel}

\usepackage{amsmath}
\usepackage{amssymb}
\usepackage[a4paper,margin=2.5cm]{geometry}
\usepackage{enumitem}
% Русские буквы в подпунктах: (а), (б), (в).
\makeatletter
\AddEnumerateCounter{\asbuk}{\@asbuk}{м}
\makeatother
\usepackage{listings}

% Чёрно-белое оформление кода: комментарии курсивом, без цвета.
\lstdefinestyle{bw}{
  language=Python,
  basicstyle=\ttfamily\small,
  commentstyle=\itshape,
  numbers=none,
  frame=single,
  columns=fullflexible,
  keepspaces=true,
  showstringspaces=false,
  breaklines=true,
  upquote=true,
}
\lstset{style=bw}

\newcommand{\pts}[1]{\textbf{(#1)}}

\begin{document}

\begin{center}
  {\Large\bfseries Лабораторная работа 3\par}
  {\large Доверительные интервалы\par}
\end{center}

\textbf{Как сдавать работу.}
\begin{enumerate}[itemsep=0.3em]
  \item Сохраните копию ноутбука на свой Диск («Файл $\to$ Сохранить копию на Диске»).
  \item Решите задания в своей копии.
  \item Перед отправкой выберите «Среда выполнения $\to$ Перезапустить и выполнить всё» и убедитесь, что выводы всех ячеек (числа, графики) видны.
  \item В конце ноутбука напишите 2--3 предложения вывода по практической части.
  \item Откройте доступ: «Поделиться $\to$ Все, у кого есть ссылка $\to$ Читатель».
  \item Пришлите преподавателю ссылку на свою копию.
\end{enumerate}

Всего за работу можно получить 10 баллов.

Обозначения: $X_i$~--- элементы случайной выборки (случайные величины), $x_i$~--- их реализация (наблюдаемые значения).

\section*{Теоретические вопросы}

\begin{enumerate}
  \item \pts{1 балл} Что такое доверительный интервал уровня доверия $1-\alpha$ для параметра $\theta$? Объясните смысл фразы «интервал, вычисленный по выборке,~--- это 95\%-й доверительный интервал для $\theta$»: что в этой конструкции случайно и что означают 95\%?

  \item \pts{2 балла} Пусть $\mathbf{X}=\left\{X_i\right\}_{i=1}^n$~--- случайная выборка из $\mathcal{N}(\mu,\sigma^2)$, где дисперсия $\sigma^2$ известна.
  \begin{enumerate}[label=(\asbuk*)]
    \item Найдите $\mathbb{E}\bar X_n$ и $\mathbb{D}\bar X_n$, где $\bar X_n=\frac1n\sum_{i=1}^n X_i$.
    \item Используя то, что сумма независимых нормальных величин имеет нормальное распределение, покажите, что $Z=\dfrac{\bar X_n-\mu}{\sigma/\sqrt{n}}\sim\mathcal{N}(0,1)$.
    \item Выведите доверительный интервал уровня $1-\alpha$ для $\mu$. Квантиль уровня $q$ распределения $\mathcal{N}(0,1)$ обозначайте $z_q$. Чему равен $z_{0.975}$?
  \end{enumerate}

  \item \pts{1 балл} Если дисперсия $\sigma^2$ неизвестна, ее заменяют оценкой
  \[
    S_n^2=\frac{1}{n-1}\sum_{i=1}^n\left(X_i-\bar X_n\right)^2.
  \]
  Известно, что статистика $T=\dfrac{\bar X_n-\mu}{S_n/\sqrt{n}}$ имеет распределение Стьюдента $t_{n-1}$ с $n-1$ степенями свободы (а не $\mathcal{N}(0,1)$).

  Выпишите доверительный интервал уровня $1-\alpha$ для $\mu$. Объясните, почему квантиль $t_{n-1,\,1-\alpha/2}$ больше $z_{1-\alpha/2}$ и что происходит с разницей между ними при росте $n$.

  \item \pts{1 балл} По формуле из вопроса 2 объясните, как меняется длина доверительного интервала
  \begin{enumerate}[label=(\asbuk*)]
    \item при росте объема выборки $n$;
    \item при росте уровня доверия $1-\alpha$;
    \item при росте $\sigma$.
  \end{enumerate}
  Во сколько раз нужно увеличить объем выборки, чтобы интервал стал вдвое уже?

  Какой объем выборки нужен, чтобы при $\sigma=2$ и $\alpha=0.05$ половина длины доверительного интервала была не больше $0.1$?

  \item \pts{1 балл} Пусть $\mathbf{X}=\left\{X_i\right\}_{i=1}^n$~--- случайная выборка из распределения Пуассона $\mathrm{Pois}(\lambda)$, и $\hat\lambda_n=\bar X_n$~--- оценка максимального правдоподобия из лабораторной работы 2, для которой $\mathbb{E}\hat\lambda_n=\lambda$ и $\mathbb{D}\hat\lambda_n=\lambda/n$.

  По центральной предельной теореме при больших $n$ оценка $\hat\lambda_n$ имеет приближенно нормальное распределение $\mathcal{N}(\lambda,\lambda/n)$.

  Заменив неизвестное $\lambda$ в дисперсии на $\hat\lambda_n$, запишите асимптотический доверительный интервал уровня $1-\alpha$ для $\lambda$.

\end{enumerate}

\section*{Практические задачи}

\begin{enumerate}
  \item \pts{1 балл} Реализуйте функцию, возвращающую границы доверительного интервала уровня $1-\alpha$ для математического ожидания $\mu$ нормального распределения при неизвестной дисперсии (вопрос 3).

  Квантиль берите из \lstinline|scipy.stats.t.ppf|.

  Выберите какие-нибудь значения $\mu$ и $\sigma^2$, сгенерируйте с помощью \lstinline|np.random.default_rng(seed)| выборку объема $n=20$ и сравните результат своей функции с \lstinline|stats.t.interval(1 - alpha, df=n - 1, loc=np.mean(x), scale=stats.sem(x))|:
\begin{lstlisting}
def ci_mean_t(x: np.ndarray, alpha: float = 0.05) -> tuple[float, float]:
    # your code
    return low, high
\end{lstlisting}

  \item \pts{1 балл} Для выбранных в задаче 1 значений $\mu$ и $\sigma^2$ сгенерируйте $K=100$ независимых выборок объема $n=20$ и для каждой постройте 95\%-й доверительный интервал функцией \lstinline|ci_mean_t|.

  Изобразите все интервалы на одном графике горизонтальными отрезками (по вертикальной оси~--- номер выборки, удобно использовать \lstinline|plt.hlines|), выделив другим цветом интервалы, не накрывшие $\mu$, и отметьте вертикальной линией истинное значение $\mu$.

  Сколько интервалов не накрыли $\mu$ и сколько вы ожидали?

  \item \pts{2 балла} Исследуйте асимптотический доверительный интервал для параметра распределения Пуассона из вопроса 5.

  Для $\lambda = 0.5,\ 2,\ 10$ и $n=5,10,20,50,100,500$ по $M=10000$ выборкам оцените долю 95\%-х интервалов $\hat\lambda_n\pm z_{1-\alpha/2}\sqrt{\hat\lambda_n/n}$, накрывших $\lambda$.

  Постройте графики доли накрытий от $n$ (по одной линии на каждое $\lambda$) и отметьте уровень $0.95$.

  При каких значениях $n\lambda$ этим интервалом, по вашим результатам, можно пользоваться?

\end{enumerate}

\section*{Вывод}

Напишите 2--3 предложения вывода по практической части: сколько из 100 интервалов не накрыли $\mu$ в задаче 2 и при каких $n\lambda$ асимптотический интервал для распределения Пуассона дает долю накрытий, близкую к заявленной.

\end{document}
